\section{Identification}\label{app:id}

Each proposition is followed by a one-sentence statement of what it means in
audit terms.

\subsection{Setup}

Fix a program and a period $t$. Let $C_t \in \{0,1\}$ indicate that the
underlying condition is present, and $F_t \in \{0,1\}$ that a finding is
recorded. Assume no false positives, so $F_t = 1$ implies $C_t = 1$. Write
$f \in \{0,1\}$ for the prior-period record $F_{t-1}$, and define
%
\[
  r_f = P(C_t = 1 \mid F_{t-1} = f), \qquad
  d_f = P(F_t = 1 \mid C_t = 1,\, F_{t-1} = f),
\]
\[
  q_f = P(F_t = 1 \mid F_{t-1} = f).
\]
%
Call $r_f$ the underlying rate, $d_f$ the detection probability and $q_f$ the
recorded rate. Write $\Rlat = r_1/r_0$, $\Robs = q_1/q_0$ and $M = d_1/d_0$.
The observable data determine $q_0$ and $q_1$ and nothing else. In
Propositions~\ref{prop:two}--\ref{prop:remed} take $q_0 > 0$ and $q_1 > 0$,
which is the empirically relevant case and is needed for the ratios to be
defined.

\begin{proposition}[Two-factor identity]\label{prop:two}
$q_f = r_f\,d_f$ for $f \in \{0,1\}$, and consequently $r_f \geq q_f$ and
$\Robs = \Rlat \, M$.
\end{proposition}

\begin{proof}
$P(F_t=1 \mid F_{t-1}=f) = P(C_t=1 \mid F_{t-1}=f)\,P(F_t=1 \mid C_t=1,F_{t-1}=f)
+ P(C_t=0 \mid F_{t-1}=f)\cdot 0 = r_f d_f$. Since $d_f \leq 1$,
$r_f = q_f/d_f \geq q_f$. Dividing $q_1 = r_1 d_1$ by $q_0 = r_0 d_0$ gives
the product form.
\end{proof}

In audit terms: the recorded finding rate is the problem rate times the catch
rate, so the problem rate is at least the recorded rate, and the recorded risk
ratio is the underlying risk ratio times the follow-up detection multiplier.

\begin{proposition}[Unrestricted bounds]\label{prop:bounds}
$\Rlat \in [\,q_1,\; 1/q_0\,]$, and both endpoints are attained.
\end{proposition}

\begin{proof}
By Proposition~\ref{prop:two}, $r_1 \in [q_1, 1]$ and $r_0 \in [q_0, 1]$, and
these are the only restrictions the data impose: for any such pair, setting
$d_f = q_f/r_f \in (0,1]$ reproduces $q_0$ and $q_1$ exactly. Hence
$\Rlat = r_1/r_0$ is minimized at $r_1 = q_1$, $r_0 = 1$, giving
$q_1$, and maximized at $r_1 = 1$, $r_0 = q_0$, giving $1/q_0$. Both
configurations are admissible, so both endpoints are attained.
\end{proof}

On the difference scale the same argument gives
$r_1 - r_0 \in [\,q_1 - 1,\; 1 - q_0\,]$, also sharp. The interval contains $1$
(and the difference interval contains zero) whenever $q_1 < 1$ and $q_0 < 1$.
In audit terms: the record tells us at least how common the problem must be in
each group, but not how much more common it is in one history group than the
other unless we also know something about detection.

\begin{proposition}[Monotone detection]\label{prop:mono}
If $d_1 \geq d_0$, then $\Rlat \in [\,q_1,\; \Robs\,]$.
Consequently:
\begin{enumerate}
\item[(i)] $\Robs < 1 \implies \Rlat < 1$;
\item[(ii)] if in addition $d_1/d_0 \leq \Mbar$ for a known $\Mbar \geq 1$,
then $\Rlat \geq \Robs/\Mbar$, so $\Robs > \Mbar \implies \Rlat > 1$.
\end{enumerate}
\end{proposition}

\begin{proof}
From Proposition~\ref{prop:two}, $\Robs = \Rlat\,(d_1/d_0)$. If $d_1/d_0 \geq
1$ then $\Rlat = \Robs/(d_1/d_0) \leq \Robs$; the lower bound is unchanged
from Proposition~\ref{prop:bounds}. Part~(i) is immediate. For part~(ii),
$\Rlat = \Robs/(d_1/d_0) \geq \Robs/\Mbar$.
\end{proof}

In audit terms: if follow-up can only raise the catch rate, the recorded risk
ratio is a ceiling on the underlying one, and a defended ceiling on how much
follow-up raises the catch rate gives a floor on the underlying ratio. The
Remediation column of Table~\ref{tab:worlds} shows why the ceiling is needed:
the recorded ratio is $4.0$ and the underlying ratio is $0.5$, reconciled by
$M = \nWorldRdetRatio{}$.

\begin{proposition}[Remediation reversal]\label{prop:remed}
Suppose $r_1 < r_0$, so the underlying condition is less likely after a
recorded finding. Then $q_1 > q_0$ if and only if $d_1/d_0 > r_0/r_1$.
\end{proposition}

\begin{proof}
$q_1 > q_0 \iff r_1 d_1 > r_0 d_0 \iff d_1/d_0 > r_0/r_1$, the last step
dividing by $r_1 d_0 > 0$.
\end{proof}

In audit terms: when the problem has been reduced, recorded findings still go
up after a prior finding whenever follow-up raises the catch rate by more than
the problem was reduced; in the Remediation column, $r_0/r_1 = 2$ and $M =
\nWorldRdetRatio{}$.

\subsection*{Remark: dropping the no-false-positive assumption}

Everything above assumes a recorded finding implies a real condition. If
findings can be recorded where none exists, the floor $r_f \geq q_f$ of
Proposition~\ref{prop:two} no longer holds, and with it the lower endpoint of
every interval above. Identification is then weaker, not stronger, unless
specificity is separately constrained. The assumption is kept because it is
the natural approximation for a formally adjudicated audit finding and because
it isolates the problem; relaxing it does not rescue anything.
Appendix~\ref{app:pred} reports the corresponding predictive results with
false positives allowed.

\subsection{Full histories}

Let $h$ denote an observed history of recorded findings up to $t-1$, of any
length, and write $q(h) = P(F_t = 1 \mid h)$. The observable law of the process
is exactly the collection $\{q(h)\}$ over histories $h$.

\begin{proposition}[Full-history nonidentification]\label{prop:full}
Let $r(\cdot)$ be any function with $r(h) \in [q(h), 1]$ for every history $h$.
There is a detection function $d(\cdot)$ taking values in $[0,1]$ such that the
model with underlying rates $r(\cdot)$ and detection $d(\cdot)$ reproduces
$\{q(h)\}$ exactly.
\end{proposition}

\begin{proof}
For a history with $r(h) > 0$, set $d(h) = q(h)/r(h)$; this lies in $[0,1]$
because $0 \leq q(h) \leq r(h)$. For a history with $r(h) = 0$, the constraint
$r(h) \geq q(h)$ forces $q(h) = 0$, and $d(h)$ may be set to any value in
$[0,1]$, say $0$. In both cases $r(h)\,d(h) = q(h)$, so by
Proposition~\ref{prop:two} the model gives $P(F_t=1 \mid h) = q(h)$ for every
$h$. The observable law is determined by $\{q(h)\}$, which is unchanged.
\end{proof}

In audit terms: however many years of findings are on file, any story about
the underlying problem that stays above the recorded rate at every point can
be made to fit the whole file exactly by adjusting the catch rate, so more
years alone do not settle which story is true. The class of $r(\cdot)$
admitted includes functions increasing in the number of prior findings,
constant functions, and decreasing functions.

\subsection*{Remark: what restricted latent-state models do recover}

Proposition~\ref{prop:full} is not a statement that longitudinal data are
uninformative about latent dynamics. Latent-structure models, hidden Markov
models among them, are generically identifiable when enough periods are
observed \citep{allman2009}. Those results assume that the observed variables
are conditionally independent given the latent state, together with
homogeneity of the transition law, a stable observation structure, and no
additional unmodeled heterogeneity. History-dependent detection violates the
conditional-independence assumption directly: $F_t$ depends on $F_{t-1}$ even
after conditioning on $C_t$, because $d$ is a function of $F_{t-1}$. That is
the precise sense in which the identification problem here is not cured by
more periods, and the precise sense in which a restricted model can cure it.
If detection is restricted, fixed as a known function, tied to covariates
observed outside the record, or measured directly, the model leaves the class
covered by Proposition~\ref{prop:full}. Leaving that class is necessary for
identification and not sufficient for it: whether a particular restricted
model is identified is a question about that model, to be settled for it
specifically. What cannot work is leaving both factors free. A class that
admits the always-present condition, $r(h) \equiv 1$ with $d(h) = q(h)$,
contains a member that reproduces any observed law, and is not identified from
the record alone.

\subsection{Detection information from outside the record}

The two results below formalize the two designs of
Section~\ref{sec:operational} that identify $M$ rather than bound it. Both
are standard; they are stated here in the paper's notation.

\begin{proposition}[Reference observation under a common measurement rule]\label{prop:ref}
Suppose a reference procedure with sensitivity $s \in (0,1]$ for the target
condition and no false positives is applied to program-years that, within
each history group $f$, represent that group, by probability sampling or
through known weights, and that the procedure does not alter the condition.
Let $s$ be common to both history groups and let $q^{\mathrm{ref}}_f$ be the
(weighted) reference finding rate after history $f$. Then $q^{\mathrm{ref}}_1
/ q^{\mathrm{ref}}_0 = \Rlat$, and $M = \Robs / \Rlat$.
\end{proposition}

\begin{proof}
Representativeness within each history group means the condition rate in the
reference observations after history $f$ is $r_f$, so $q^{\mathrm{ref}}_f =
r_f\, s$ by the argument of Proposition~\ref{prop:two}. The ratio is $r_1 s /
r_0 s = \Rlat$ because $s$ is common. Then $M = \Robs/\Rlat$ by
Proposition~\ref{prop:two}.
\end{proof}

In audit terms: a reference sample that represents each history group, worked
under the same rule in both, reveals how much more common the problem is
after a prior finding even though the reference procedure misses some of it,
because what it misses cancels in the ratio; dividing the targeted record's
ratio by that number gives the follow-up detection multiplier. The sample may
be stratified on history; what it may not do is have a procedure whose
sensitivity differs between the groups. If $s_1 \ne s_0$, the reference ratio
equals $\Rlat \cdot (s_1/s_0)$ and the cancellation fails.

\begin{proposition}[Two conditionally independent procedures]\label{prop:twoproc}
Within one history group, apply two procedures A and B to the same
program-years, with detection probabilities $a$ and $b$ given $C = 1$ that
are constant across the program-years pooled, no false positives, and
$A \perp B \mid C$. Let $P_A$, $P_B$ and $P_{AB}$ be the
probabilities that A finds a condition, that B finds one, and that both do.
Then $a = P_{AB}/P_B$, $b = P_{AB}/P_A$ and $r = P_A P_B / P_{AB}$.
\end{proposition}

\begin{proof}
$P_A = r a$, $P_B = r b$ and, by conditional independence, $P_{AB} = r a b$.
Then $P_{AB}/P_B = a$, $P_{AB}/P_A = b$, and $P_A P_B / P_{AB} = r$.
\end{proof}

In audit terms: among the problems the second procedure finds, the share the
first also found is the first procedure's catch rate, and once both catch
rates are known the number of problems both missed follows; this is the
occupancy-and-detection logic of \citet{mackenzie2002}, and it fails when the
two procedures miss the same cases for the same reason, because then $P_{AB}$
exceeds $r a b$ and both catch rates are overstated, or when sensitivity
varies across the pooled program-years so that no single $a$ or $b$ exists. Applied in each history
group it yields $d_0$, $d_1$ and $M$ for whichever procedure the targeted
audit uses.

\subsection{Serial dependence generated by observation alone}

\begin{proposition}[Observation-only dynamics]\label{prop:markov}
Suppose $C_t$ is independent across $t$ with $P(C_t=1) = p$, and detection
depends on the record through $d_{F_{t-1}}$ alone. Then $\{F_t\}$ is a two-state
Markov chain with $q_0 = p\,d_0$ and $q_1 = p\,d_1$; its autocorrelation is
$\rho_k = \lambda^k$ with $\lambda = q_1 - q_0 = p\,(d_1 - d_0)$; and
$\rho_1 \leq p$.
\end{proposition}

\begin{proof}
$P(F_t=1 \mid F_{t-1}=f) = p\,d_f$ by Proposition~\ref{prop:two} and the
independence of $C_t$ from the past. A two-state chain with transition
probabilities $q_0$ and $q_1$ has lag-one autocorrelation $q_1 - q_0$ and
$\rho_k = (q_1-q_0)^k$. Finally $\lambda = p(d_1-d_0) \leq p$ since
$d_1 - d_0 \leq 1$.
\end{proof}

In audit terms: even when the problem comes and goes at random with no
memory at all, a follow-up rule makes the finding record look like it has
memory, and the recorded risk ratio ($M$ here, since $\Rlat = 1$) says
nothing about how strong that apparent memory is. In the exact case with
$p = \nObsOnlyP{}$, $d_0 = \nObsOnlyDzero{}$ and $d_1 = \nObsOnlyDone{}$, the
recorded risk ratio is \nObsOnlyRR{} while the lag-one autocorrelation of the
record is \nObsOnlyRho{}. Neither the magnitude of the recorded elevation nor
the presence of serial dependence in the record discriminates between a
persisting condition and a persisting audit response.

\subsection*{Remark: stable heterogeneity is a fourth channel}

Persistence, observation feedback and remediation are not the only ways a
finding record can look autocorrelated. Entities may simply differ: some
auditees are durably more prone to deficiencies than others, and that alone
produces association between one year's record and the next. Adding such
heterogeneity to the multitype setting of Table~\ref{tab:multitype} raises
findings in requirement areas that had never had one, so it is a candidate
explanation for the pattern that generalized scrutiny also produces. At the
prespecified parameters it does not reproduce that pattern in magnitude: the
raw new-area risk ratio is \nNullCategoryRR{} with neither channel (below one,
because of the opportunity-set artifact of Section~\ref{sec:fac}),
\nFrailtyRR{} with heterogeneity alone, and \nScrutinyRR{} with generalized
scrutiny. This is a negative result at one set of parameters, not a
demonstration that the two channels are distinguishable in general.
